\answer{OA28}
The total weight is $\sum_{n=0}^\infty2^{-n}=2$. Thus the series converges at every $x$, with a tail bound that is uniform in $x$.

First let $x$ be irrational. Given $\varepsilon>0$, choose $N$ such that $\sum_{n>N}2^{-n}<\varepsilon$. None of the finitely many points $r_0,\ldots,r_N$ equals $x$. Hence some $\delta>0$ ensures that, whenever $|y-x|<\delta$, each of these rational points is on the same side of $y$ as it is of $x$. The first $N+1$ terms remain unchanged, so
\[
 |f(y)-f(x)|\le\sum_{n>N}2^{-n}<\varepsilon.
\]
The function is continuous at every irrational point.

Now take $x=r_k$ and separate its term:
\[
 f(y)=2^{-k}\mathbf1_{\{y>r_k\}}+g(y).
\]
The same finite-prefix and tail argument shows that $g$ is continuous at $r_k$, because its prefix contains no point equal to $r_k$. The strict inequality in the definition gives $f(r_k)=g(r_k)$. Therefore
\[
 \boxed{\lim_{y\uparrow r_k}f(y)=f(r_k),\qquad
 \lim_{y\downarrow r_k}f(y)=f(r_k)+2^{-k}.}
\]
Each rational point has an upward jump of size $2^{-k}$ and is left-continuous but not continuous. The continuity points are exactly the irrationals. Dense jump locations cause no contradiction: their tail weights can be made uniformly small.
